% Source r6 block. Main proof is retained unchanged in its new position.
\section{Excluding a bounded negative part}\label{sec:bounded}

We now prove Theorem~\ref{res:bounded}, stated in the introduction.
The argument uses the defect identity, absorption, gcd stabilisation and the
bounded-rise barrier.  The ordinary canonical-tail transfer supplies
integrality and normalised vanishing under the original hypotheses; only
the eventual lower bound on $E_n$ is additional.

On this common exact-orbit class, bounded negative error and finite
normalised negative mass are each equivalent to eventual zero error.
They are different quantities to estimate and have different proofs, not
incomparable endpoint classes.  The scalar finite-mass lemma below also
works without the denominator dynamics or normalised vanishing.

Hypotheses~(5) and~(6) of Theorem~\ref{res:bounded} pull against each other.  By
Proposition~\ref{res:update}, hypothesis~(5) caps the rise of the tail state at
$B$ per step, while hypothesis~(6), once the error is nowhere zero, forces the
tail state to diverge.  Proposition~\ref{res:gcdstab} makes the orbit reduced,
and Theorem~\ref{res:barrier} then converts the tension into a contradiction:
an integer sequence tending to infinity with upward steps of size at most $B$
cannot remain coprime to the pairwise coprime moduli a reduced tail carries.
The
complementary case, in which the error is eventually positive, is the descent
of Theorem~\ref{res:descent}.

\begin{proof}[Proof of Theorem~\ref{res:bounded}]
Suppose not.  We first remove every late zero: by Theorem~\ref{res:absorb} and
its contrapositive form, $E$ is nowhere zero beyond the centring threshold.
Two cases remain.  If $E$ is negative cofinally,
then~(5) bounds those magnitudes, so Proposition~\ref{res:gcdstab} makes the
tail gcd stabilise and the orbit may be taken reduced past that point.
Hypothesis~(5) with
Proposition~\ref{res:update} gives $C_{n+1}\le C_n+B$, a bounded rise;
hypothesis~(6) with $E$ nowhere zero gives $C_n\to\infty$.  A reduced exact
tail with a divergent numerator and bounded rise is excluded by
Theorem~\ref{res:barrier}.  It remains to treat the case that $E$ is eventually
positive, and there the descent
of Theorem~\ref{res:descent} applies and forces $E$ to vanish, contradicting
nowhere-vanishing.
\end{proof}

\noindent Formalised as the
\lword{Erdos243/ReciprocalTailRigidity.lean}{2265}{boundedNegativePart_eventually_zero}{bounded-negative-part
rigidity}, with the eventual form, in which the centring and the bound need
only hold from some index, as the
\lword{Erdos243/ReciprocalTailRigidity.lean}{2352}{eventuallyBoundedNegativePart_eventually_zero}{eventual
bounded-negative-part rigidity}.  The two auxiliary results it uses are the
\lword{Erdos243/ReciprocalTailRigidity.lean}{1733}{tailState_tendsto_atTop_of_nonzero_normalizedVanishes}{divergence
from normalised vanishing} and the
\lword{Erdos243/ReciprocalTailRigidity.lean}{1748}{no_cofinallyBoundedNegative_of_normalizedVanishes}{exclusion
of cofinally bounded negative magnitudes}, the latter over the
\lword{Erdos243/ReciprocalTailRigidity.lean}{1655}{no_cofinallyBoundedNegative_of_tailTendsto}{tail-divergence
form}.

\begin{corollary}\label{res:cor}
Under the hypotheses of Theorem~\ref{res:bounded} the multipliers satisfy
$a_{n+1}=a_n^{2}-a_n+1$ for all sufficiently large $n$.
(The extra $C_{n+1}\ne0$ sometimes written beside Theorem~\ref{res:eventual}
is redundant: hypothesis~(1) already gives $C_n>0$.)
\end{corollary}

\begin{proof}
Theorem~\ref{res:bounded} then Theorem~\ref{res:eventual}.
\end{proof}

Hypotheses~(1)--(3) are the state system, and~(6) is the division-free form of
$|E_n|/C_n\to0$.  Hypothesis~(4) is logically redundant: hypothesis~(6) with
$K=1$ gives $|E_n|<C_n$ eventually.  It is retained because the displayed
statement tracks the linked Lean declaration; shifting to the resulting
threshold gives the same conclusion without an independent centring input.
The ordinary implication from rationality and quadratic growth to integrality
of $C_n$ and $E_n/C_n\to0$ is settled mathematics (Koizumi's canonical-tail
transfer; equivalently the two-term tail estimate).  Lean does not yet
transport that bridge, so Theorem~\ref{res:bounded} remains a conditional
theorem about the state system.  Combined with the ordinary bridge,
Corollary~\ref{res:cor} settles the bounded-negative canonical case of
Problem~\ref{res:problem}.  The state implication and its ordinary original-coordinate consequence are
proved theorems with displayed hypotheses.  The missing formal bridge does
not make the ordinary implication an open mathematical question.

The conditionality can nevertheless be located precisely, because the missing
inputs are supplied elsewhere.  With the results of~\cite{koizumi2025},
Theorem~\ref{res:bounded} becomes a conditional theorem about
Problem~\ref{res:problem} itself rather than about the state system alone.  Let
$(a_n)$ satisfy the hypotheses of Problem~\ref{res:problem}.  After deleting a
finite prefix, Corollary~3 of~\cite[p.~9]{koizumi2025} makes the sequence the
pseudo-greedy expansion of its own reciprocal sum, and Lemma~4 there supplies
the integers $c_n,d_n,e_n$ of Section~\ref{sec:problem}.  Hypothesis~(1) then
holds: $c_n$ is a positive integer by Lemma~4(1), and $a_n>1$ after a further
finite shift, because summability of $\sum1/a_n$ forces $a_n\to\infty$.
Hypothesis~(2) is the pair of recurrences $c_{n+1}=a_nc_n-d_n$ and
$d_{n+1}=a_nd_n$ of Lemma~4(2), and hypothesis~(3) is that lemma's
$a_n=(d_n-e_n)/c_n+1$ read as $e_n=d_n-(a_n-1)c_n$, which is the map $\ctr$.
The centring range $-c_n/2\le e_n<c_n/2$ in the same lemma gives
hypothesis~(4) at every index, which is stronger than the eventual form used
here.  Hypothesis~(6) is the vanishing
of the gap sequence in Corollary~3, since $\varepsilon_n=e_n/c_n$.  The sole
hypothesis not supplied is~(5), the eventual bound on the negative part;
nothing in~\cite{koizumi2025} yields it and nothing here does either, so
Problem~\#243 stays open.
